\documentclass[11pt]{article}
\usepackage[a4paper,margin=1in]{geometry}
\usepackage{amsmath,amssymb,amsthm,mathtools}
\usepackage{graphicx}
\usepackage[hidelinks]{hyperref}
\usepackage{microtype}

\newtheorem{theorem}{Theorem}
\newtheorem{lemma}{Lemma}
\newcommand{\R}{\mathbb R}
\newcommand{\Leb}{\mathcal L}
\newcommand{\Mloc}{\mathcal M_{\mathrm{loc}}}
\newcommand{\Dom}{\operatorname{Dom}}

\title{Laplacian closability: the finite-exponent theorem\\
and the strong-$L^\infty$ counterexample}
\author{Full resolution of the expectation in arXiv:2311.08058}
\date{21 July 2026}

\begin{document}
\maketitle

\begin{center}
\fbox{\parbox{0.91\textwidth}{\textbf{Status.}
Candidate counterexample and full resolution, likely valid, awaiting human
review.  The expected necessary condition holds for every finite output
exponent in all dimensions, but fails at the strong $L^\infty$ endpoint.}}
\end{center}

\section{Source expectation}

Alberti, Bate, and Marchese study closability of differential operators with
respect to Radon measures in \cite{ABM}.  In concluding Remark~1.7(iv), page~4,
they suggest that closability of the Laplacian from $L^p(\mu)$ to $L^q(\mu)$
should force $\mu\ll\Leb^d$.

\begin{center}
\includegraphics[width=\textwidth]{figures/open_problem_crop.png}
\end{center}

\noindent Here closability is sequential: for a linear operator it is enough
to require that
\[
 u_n\longrightarrow0\text{ in the input space},\qquad
 \Delta u_n\longrightarrow w\text{ in the output space}
 \quad\Longrightarrow\quad w=0.
\]

\section{Idea of the proof}

For finite output exponent, closability forces the Banach adjoint to have a
dense domain.  Every adjoint test function $\varphi$ then gives a measure
$\nu=\varphi\mu$ whose distributional Laplacian is also a Radon measure.
The fundamental solution of the Laplacian shows that such a $\nu$ must have a
Lebesgue density.  Closedness plus density of the adjoint domain then forces
every compact restriction of $\mu$ to have a density.

The $L^\infty$ endpoint behaves differently.  For
$\mu=\Leb^d+\delta_0$, the $L^\infty(\mu)$ norm of a continuous function is
its ordinary supremum norm.  Thus convergence of the Laplacians is uniform.
Input convergence makes the original functions converge to zero as
distributions, so the uniform limit of their Laplacians must be zero.  The
operator is closable although $\mu$ has an atom.

\section{Full dichotomy}

\begin{theorem}\label{thm:main}
Let $d\ge1$, let $\mu$ be a positive Radon measure on $\R^d$, and initially
define $\Delta$ on $C_c^2(\R^d)$.
\begin{enumerate}
\item If $1\le p\le\infty$ and $1\le q<\infty$, closability of
\[
 \Delta:L^p(\mu)\supset C_c^2(\R^d)\longrightarrow L^q(\mu)
\]
forces $\mu\ll\Leb^d$.
\item The assertion is false for $q=\infty$.  For
\[
 \mu=\Leb^d+\delta_0
\]
the measure $\mu$ is not absolutely continuous, but $\Delta$ is closable from
$L^p(\mu)$ to $L^\infty(\mu)$ for every $1\le p\le\infty$.
\end{enumerate}
\end{theorem}

We use two elementary lemmas.  The first is the standard adjoint criterion,
included for clarity.

\begin{lemma}[Adjoint density]\label{lem:adjoint}
Let $X,Y$ be Banach spaces and let $T:\Dom(T)\subset X\to Y$ be densely
defined.  If $T$ is closable, then $\Dom(T^*)$ is weak-star dense in $Y^*$.
If $Y$ is reflexive, this linear subspace is norm dense.
\end{lemma}

\begin{proof}
Let $G\subset X\times Y$ be the graph.  If $y\in Y$ annihilates
$\Dom(T^*)$, then $(0,y)$ annihilates every element of $G^\perp$: indeed,
the second coordinate of an element of $G^\perp$ belongs to $\Dom(T^*)$.
The bipolar theorem for linear subspaces gives $(0,y)\in\overline G$.
Since $\overline G$ is a graph, $y=0$.  Thus the preannihilator of
$\Dom(T^*)$ is zero, which is exactly weak-star density.  In a reflexive
space, a linear subspace with weakly dense closure is norm dense by
Hahn--Banach.
\end{proof}

\begin{lemma}[Measure with measure Laplacian]\label{lem:elliptic}
If $\nu\in\Mloc(\R^d)$ and $\Delta\nu\in\Mloc(\R^d)$ as distributions, then
$\nu\ll\Leb^d$.
\end{lemma}

\begin{proof}
Write $\eta=\Delta\nu$.  Fix a ball $B$ and choose
$\chi\in C_c^\infty(\R^d)$ equal to one on a neighborhood of $B$.  Let $E$
be the fundamental solution of the Laplacian and set
\[
 v:=E*(\chi\eta).
\]
The function $E$ is locally integrable in every dimension, so Fubini's theorem
shows that $v\in L^1_{\mathrm{loc}}$.  On $B$,
$\Delta(\nu-v)=0$.  Weyl's lemma for harmonic distributions makes
$\nu-v$ smooth on $B$.  Hence $\nu$ is represented by an $L^1_{\mathrm{loc}}$
function on $B$.  Since $B$ was arbitrary, $\nu\ll\Leb^d$.
\end{proof}

\begin{proof}[Proof of Theorem~\ref{thm:main}, finite $q$]
It suffices to restrict the operator to $C_c^\infty(\R^d)$.
If two smooth functions representing the same $L^p(\mu)$ class already have
different Laplacians in $L^q(\mu)$, the graph is not a graph and closability
fails immediately; thus under the closability assumption the restriction is a
well-defined operator.

First suppose $p<\infty$.  Put $X=L^p(\mu)$ and $Y=L^q(\mu)$.  Smooth
compactly supported functions are dense in $X$.  If the graph were closable,
Lemma~\ref{lem:adjoint} would make $\Dom(\Delta^*)$ weak-star dense in
$Y^*=L^{q'}(\mu)$; for $q>1$ this is norm density.

For $\varphi\in\Dom(\Delta^*)$ there is $g\in L^{p'}(\mu)$ such that
\[
 \int_{\R^d}\varphi\,\Delta u\,d\mu
 =\int_{\R^d}g u\,d\mu
 \qquad (u\in C_c^\infty(\R^d)).
\]
Thus, for the locally finite signed measure $\nu=\varphi\mu$,
\[
 \Delta\nu=g\mu\in\Mloc(\R^d).
\]
Lemma~\ref{lem:elliptic} gives $\varphi\mu\ll\Leb^d$.

Now suppose $p=\infty$.  Define the same operator on the dense subspace
$C_c^\infty(\R^d)$ of $X=C_0(\R^d)$.  If $u_n\to0$ uniformly, then
$u_n\to0$ in $L^\infty(\mu)$, so the assumed closability makes this
$C_0$-operator closable.  For an adjoint-domain element $\varphi$, the Riesz
representation theorem supplies a finite signed measure $\eta$ with
\[
 \int\varphi\,\Delta u\,d\mu=\int u\,d\eta.
\]
Hence $\Delta(\varphi\mu)=\eta$, and Lemma~\ref{lem:elliptic} again yields
$\varphi\mu\ll\Leb^d$.

In either case, let
\[
 \mathcal A:=\{\varphi\in L^{q'}(\mu):\varphi\mu\ll\Leb^d\}.
\]
For $q>1$, $\mathcal A$ is norm closed: norm convergence in $L^{q'}(\mu)$
implies total-variation convergence of the associated measures on every
compact set.  For $q=1$, $\mathcal A\subset L^\infty(\mu)$ is weak-star
closed: for every Lebesgue-null Borel set $N$ of finite $\mu$-measure, the
condition $\int_N\varphi\,d\mu=0$ is weak-star closed, and one intersects these
conditions over all such $N$ and their Borel subsets.

The adjoint domain lies in $\mathcal A$ and is dense in the relevant topology,
so $\mathcal A=L^{q'}(\mu)$.  For each compact $K$, the function
$\mathbf1_K$ belongs to $L^{q'}(\mu)$; therefore
$\mu\mathbin{\vrule height 1.4ex depth -0.3ex width 0.07ex\vrule height 0.07ex depth -0.02ex width 0.7ex}K\ll\Leb^d$.
An exhaustion by compact balls gives $\mu\ll\Leb^d$.
\end{proof}

\begin{proof}[Proof of Theorem~\ref{thm:main}, $q=\infty$]
Let $\mu=\Leb^d+\delta_0$.  This is a positive Radon measure and is not
absolutely continuous.  For every continuous compactly supported function
$f$,
\begin{equation}\label{eq:sup}
 \|f\|_{L^\infty(\mu)}=\|f\|_\infty,
\end{equation}
because the Lebesgue essential supremum of a continuous function is its
ordinary supremum.

Suppose $u_n\in C_c^2(\R^d)$, $u_n\to0$ in $L^p(\mu)$, and
$\Delta u_n\to w$ in $L^\infty(\mu)$.  If $p<\infty$, then
$u_n\to0$ in $L^p(\Leb^d)$ and hence in distributions.  If $p=\infty$,
\eqref{eq:sup} applied to $u_n$ gives uniform convergence to zero, hence again
distributional convergence.

By \eqref{eq:sup}, the continuous functions $\Delta u_n$ are uniformly
Cauchy.  They converge uniformly to a continuous function $f$ representing
$w$.  Distributionally, however,
\[
 \Delta u_n\longrightarrow0
\]
because $u_n\to0$ in distributions.  Uniform convergence also gives
$\Delta u_n\to f$ in distributions, so $f=0$ and therefore $w=0$.
This proves closability for every $p$.
\end{proof}

\section{Consequences and limitations}

The theorem completely resolves the necessary-condition expectation as
stated across the output exponent range.  The finite-exponent result includes
$q=1$ and the input endpoint $p=\infty$.  The strong-$L^\infty$ endpoint is
not merely outside the proof: the implication is false there.

No converse is asserted.  Even for absolutely continuous measures,
closability may depend on the density and on the exponents.  The theorem also
does not concern weak-star convergence in $L^\infty$, which is a different
output topology.

\section{Novelty check and review recommendation}

The local run indexes were searched for arXiv:2311.08058, Laplacian
closability, measure Laplacians, and the strong-$L^\infty$ endpoint.  Bounded
arXiv searches through 21 July 2026 used exact and close variants of
``Laplacian closable Radon measure'', ``absolute continuity'', and
``$L^\infty$ closability''.  They found the source paper but no later exact
answer.  This is not an exhaustive bibliographic claim.

Human review should focus on Lemma~\ref{lem:elliptic}, the weak-star closure
argument at $q=1$, and the fact that full-support Lebesgue measure turns the
output $L^\infty$ norm into the uniform norm.  Subject to those checks, this
packet should replace the earlier one-dimensional partial result.

\begin{thebibliography}{9}
\bibitem{ABM}
G. Alberti, D. Bate, and A. Marchese,
\emph{On the closability of differential operators},
J. Funct. Anal. 289 (2025), 111029;
\href{https://arxiv.org/abs/2311.08058}{arXiv:2311.08058v2}.

\bibitem{Evans}
L. C. Evans,
\emph{Partial Differential Equations}, 2nd ed.,
Graduate Studies in Mathematics 19, American Mathematical Society, 2010.
\end{thebibliography}

\end{document}
